% from ega2-1-fr.tex
\oldpage[II]{116}

des $f^{-1}(y')$, où $y'\neq y$ est dans $g^{-1}(z)$~; comme $x$ est isolé
dans $f^{-1}(y)$, il est isolé dans
$h^{-1}(z)=f^{-1}(g^{-1}(z))$. Pour prouver (iii), on peut se limiter au cas
où $Y=S$ (\hyperref[I.3.3.11-fr]{I, 3.3.11})~; on remarque encore que tout
d'abord $f'=f_{(S')}$ est de type fini
(\hyperref[I.6.3.4-fr]{I, 6.3.4, (iii)})~; d'autre part, si
$x'\in X'=X_{(S')}$ et si on pose $y'=f'(x')$ et $y=g(y')$,
$f'^{-1}(y')$ s'identifie à
$f^{-1}(y)\otimes_{k(y)}k(y')$ (\hyperref[I.3.6.5-fr]{I, 3.6.5})~; comme
$f^{-1}(y)$ est de rang fini sur $k(y)$ par hypothèse, $f'^{-1}(y')$ est de
rang fini sur $k(y')$, donc discret. Les assertions (iv), (v), (vi) se
déduisent de (i), (ii), (iii) par la méthode générale
(\hyperref[I.5.5.12-fr]{I, 5.5.12}), sauf lorsque les hypothèses dans (v) sont
autres que l'hypothèse «~$g$ séparé~»~; pour traiter ce cas, on remarque
d'abord que si $x$ est isolé dans
$f^{-1}(g^{-1}(g(f(x))))$, il l'est a fortiori dans $f^{-1}(f(x))$~; le fait
que $f$ est de type fini résulte alors de
(\hyperref[I.6.3.6-fr]{I, 6.3.6}).

\begin{proposition}[6.2.5]
\label{II.6.2.5-fr}
Soient $A$ un anneau local noethérien complet, $X$ un $A$-schéma localement
de type fini, $x$ un point de $X$ au-dessus du point fermé $y$ de
$Y=\operatorname{Spec}(A)$, et supposons $x$ isolé dans sa fibre
$f^{-1}(y)$ ($f$ étant le morphisme structural $X\to Y$). Alors
$\mathcal{O}_x$ est un $A$-module de type fini et $X$ est $Y$-isomorphe à la
somme (I, 3.1) de
$X'=\operatorname{Spec}(\mathcal{O}_x)$ (qui est un $Y$-schéma fini) et d'un
$A$-schéma $X''$.
\end{proposition}

Il résulte de (\hyperref[II.6.2.1-fr]{6.2.1}) que $\mathcal{O}_x$ est un
$A$-module quasi-fini. Comme $\mathcal{O}_x$ est noethérien
(\hyperref[I.6.3.7-fr]{I, 6.3.7}) et l'homomorphisme
$A\to\mathcal{O}_x$ local, l'hypothèse que $A$ est complet entraîne que
$\mathcal{O}_x$ est un $A$-module de type fini
(\hyperref[0.7.4.3-fr]{0, 7.4.3}). Soit
$X'=\operatorname{Spec}(\mathcal{O}_x)$ le schéma local de $X$ au point $x$
(\hyperref[I.2.4.1-fr]{I, 2.4.1}), et soit $g:X'\to X$ le morphisme
canonique. Comme le composé $X'\xrightarrow{g}X\xrightarrow{f}Y$ est fini
(\hyperref[II.6.1.1-fr]{6.1.1}) et que $f$ est séparé, $g$ est fini
(\hyperref[II.6.1.5-fr]{6.1.5, (v)}), donc $g(X')$ est fermé dans $X$
(\hyperref[II.6.1.10-fr]{6.1.10})~; d'autre part, comme $g$ est de type fini,
$g$ est un isomorphisme local au point fermé $x'$ de $X'$ en vertu de la
définition de $g$ et de (\hyperref[I.6.5.4-fr]{I, 6.5.4})~; mais comme $X'$
est le seul voisinage ouvert de $x'$, cela signifie que $g$ est une immersion
ouverte, donc $g(X')$ est aussi ouvert dans $X$, ce qui achève la
démonstration.

\begin{corollary}[6.2.6]
\label{II.6.2.6-fr}
Soient $A$ un anneau local noethérien complet, $Y=\operatorname{Spec}(A)$,
$f:X\to Y$ un morphisme séparé et quasi-fini. Alors $X$ est $Y$-isomorphe à
une somme $X'\amalg X''$, où $X'$ est un $Y$-schéma fini et $X''$ un
$Y$-schéma quasi-fini tel que, si $y$ est le point fermé de $Y$,
$X''\cap f^{-1}(y)=\varnothing$.
\end{corollary}

En effet, la fibre $f^{-1}(y)$ est finie et discrète par hypothèse, et le
corollaire résulte donc, par récurrence sur le nombre de points de cette
fibre, de (\hyperref[II.6.2.5-fr]{6.2.5}).

\begin{remark}[6.2.7]
\label{II.6.2.7-fr}
Au chapitre V, nous verrons que si $Y$ est localement noethérien, un morphisme
quasi-fini séparé $X\to Y$ est nécessairement quasi-affine.
\end{remark}

\subsection{Fermeture intégrale d'un préschéma.}
\label{subsection:II.6.3-fr}

\begin{proposition}[6.3.1]
\label{II.6.3.1-fr}
Soient $(X,\mathcal{A})$ un espace annelé, $\mathcal{B}$ une
$\mathcal{A}$-Algèbre (commutative), $f$ une section de $\mathcal{B}$
au-dessus de $X$. Les propriétés suivantes sont équivalentes~:
\begin{enumerate}
\item[a)] Le sous-$\mathcal{A}$-Module de $\mathcal{B}$ engendré par les
$f^n$ pour $n\geq 0$ (\hyperref[0.5.1.1-fr]{0, 5.1.1}) est de type fini.
\item[b)] Il existe une sous-$\mathcal{A}$-Algèbre $\mathcal{C}$ de
$\mathcal{B}$, qui est un $\mathcal{A}$-Module de type fini, telle que
$f\in\Gamma(X,\mathcal{C})$.
\item[c)] Pour tout $x\in X$, $f_x$ est entier sur la fibre
$\mathcal{A}_x$.
\end{enumerate}
\end{proposition}

